\documentclass[../main.tex]{subfiles}
\begin{document}
%==========================================TIÊU ĐỀ TẬP========================================
\begin{table}[h]
    \begin{adjustwidth}{-1cm}{}
        \begin{tabular}{>{\centering\arraybackslash}p{6cm}|>{\raggedright\arraybackslash}p{12.5cm}}
            \multirow{5}{6cm}
            % Cột bên trái: ÉPISODE và số tập
            {\\[-40pt]
                \flaregothic\fontsize{20pt}{20pt}\selectfont \centering
                \color{eptype}HISTOIRE\color{black}\\
                \freshbot\fontsize{72pt}{72pt}\selectfont
                \color{epnum}3\\[6pt]
                \cafeteria\fontsize{20pt}{20pt}\selectfont\color{epnum}(D-3)\color{black}
                \\[18pt]
                % \flaregothic\fontsize{14pt}{14pt}\selectfont
                % \color{epattr}BIENVENUE, \uline{\color{Banpire} Mel} !
                % \flaregothic\fontsize{18pt}{18pt}\selectfont
                % \color{epattr}ÉPISODE PERDU
                \color{black}
            }
             &
            % Dòng thứ nhất: tên tập tiếng Nhật
            {\fotskip\fontsize{18pt}{18pt}\selectfont \ruby{火}{ひ}\ruby{川}{かわ}\ruby{寮}{りょう}\ruby{カ}{か}\ruby{オ}{お}\ruby{ス}{す}\ruby{拡}{かく}\ruby{張}{ちょう}}
            \\[6pt]
            % Dòng thứ hai: tên tập tiếng Anh
             & {\cafeteria\fontsize{18pt}{18pt}\selectfont Expand the House!}                                                \\[4pt]
            % Dòng thứ ba: tên tập tiếng Việt
             & {\vnmsans\fontsize{18pt}{18pt}\selectfont Xây căn phòng mới}                                                  \\[8pt]
            % Dòng thứ tư: Nhân vật xuất hiện
             & {\caviar{\fontsize{14pt}{14pt}\selectfont Nhân vật: \emp{kuon1} \emp{ryuuhou1} \emp{aloe}}}
            \\[8pt]
        \end{tabular}
    \end{adjustwidth}
\end{table}
%=========================================TRUYỆN CHÍNH========================================
\color{black}

\txtbx{
{\color{vol} \uline{Tóm tắt:}} YAGOO gọi điện bất ngờ đề nghị Hikawa mở rộng thêm một tầng để đón một vị khách bí ẩn. Kuon, Ryujin, Ryuuhou và Aloe bắt tay vào dự án với tay nghề nghiệp dư: Aloe vẽ bản thiết kế lâu đài ma thuật, máy xây dựng tự động dựng ra một căn phòng nghiêng 45 độ với cửa cao ba mét. Hijaki cười ngất và nhóm sinh viên kiến trúc phải đến đo đạc lại toàn bộ nhà. Sau một tuần chỉnh sửa, căn phòng mới hoàn thiện đúng chuẩn, chờ đợi cho vị khách tiếp theo sẽ sinh sống tại nơi đây\dots}

\pov{Hikawa Kuon}

Chiều thứ Sáu cuối tháng Năm, một ngày hiếm hoi tôi được ở nhà nghỉ ngơi. Trong khi năm thành viên của Thế hệ thứ Nhất đang nhiệt tình tham gia buổi livestream ``From 1st'', tôi lại tận hưởng khoảng thời gian yên tĩnh riêng, không phải bận tâm về công việc hay đống trách nhiệm chồng chất. Tiếng nhạc từ buổi phát trực tuyến phát ra từ điện thoại, nhưng tôi chỉ để nó làm nền, chẳng thực sự để ý vào nội dung.

Bố tôi cũng kết thúc công việc sớm nên ở nhà cùng tôi, mẹ vẫn đang bận rộn ở kho hàng, còn em gái tôi thì vẫn phải đến trường như mọi ngày.

Khoảnh khắc bình yên này gợi lại kỷ niệm mười năm trước, khi chỉ có hai bố con tôi mới dọn đến đây. Lúc đó mẹ và em gái chưa thể chuyển đến sống chung vì vài lý do cá nhân, nhưng theo dòng thời gian, cả gia đình cuối cùng cũng được đoàn tụ dưới một mái nhà.

Và rồi, Aloe - nàng Succubus - trở thành người đầu tiên ngoài gia đình Hikawa chuyển đến sống chung. Sau khi rời khỏi \hll, cô gái mười lăm tuổi ấy đã dần trở thành một phần không thể thiếu trong cuộc sống của mọi người, như một làn gió mới thổi vào ngôi nhà vốn chỉ có sự yên ả.

Đó là những câu chuyện của quá khứ, còn hiện tại thì sao?

Bây giờ, tôi đang trải dài trên ghế, cầm chiếc điện thoại máy game tôi ``mượn'' được từ kệ của nàng phù thủy sau một sự kiện khó kể lể\footnote{{\flaregothic ÉPISODE 16}, quyển số Một.} để giải trí. Còn bố tôi thì dựa lưng vào ghế, mắt dán vào màn hình TV xem một bộ phim cũ ông từng mê mẩn hồi trẻ, thỉnh thoảng lại bật cười một mình với nội dung trên phim.

Một buổi chiều yên bình, tưởng chừng sẽ trôi qua nhẹ nhàng như những ngày trước khi tôi bước chân vào công ty\dots

Nhưng tiếng chuông điện thoại đột ngột vang lên, phá tan bầu không khí tĩnh lặng.

Tôi liếc nhìn màn hình, và một cái tên quá đỗi quen thuộc hiện lên: YAGOO.

Tôi thở dài. Nếu vị CEO này gọi điện, chắc chắn sẽ không phải chuyện đơn giản. Có lẽ ông lại muốn nhờ tôi làm gì đó, có thể là một dự án mới, hay một vấn đề nào đó mà tôi không thể giải quyết một mình.

Tôi nhấn nút trả lời, đã chuẩn bị tinh thần cho mọi rắc rối có thể xảy ra: ``Chú Tanigo-san đấy ạ?''

``À, Hikawa-san! Tôi có một đề xuất này, nhà các cậu có nghĩ đến việc mở rộng thêm một tầng nữa không?''

Tôi hơi nheo mắt trước đề xuất kỳ lạ, vô thức đặt điện thoại xuống bàn để nghe rõ hơn: ``Không ạ. Sao lại có ý đó vậy chú?''

``Tôi có vài ý tưởng cho một dự án mới, và thấy rằng nếu nhà Hikawa thêm một tầng, đó sẽ là nơi hoàn hảo cho dự án ấy.''

``Dự án gì vậy chú?''

``Để tôi nói sau. Nếu nhà cậu đồng ý, ta sẽ cử một nhóm kiến trúc sư qua khảo sát và thiết kế. Còn nếu không thì cũng không sao cả.''

Tôi im lặng một lúc, vẫn chưa hoàn toàn tin vào những gì ông nói. Giọng ông trông càng hào hứng, tôi càng cảm thấy phải đề phòng: ``Cụ thể hơn được không ạ?''

Ông ấy trả lời với giọng hơi vội vàng, như đang cố che giấu điều gì đó: ``À thì\dots\ sẽ có người cần chỗ ở thôi! Thôi cứ xây đi, sau này cậu sẽ hiểu!''

Càng nghe ông nói, tôi càng nhíu mày, cảm thấy như mình đang bị kéo vào một kế hoạch lớn mà tôi chưa thể nhìn rõ toàn bộ bức tranh: ``\dots Chú cứ úp mở như thế này làm sao cháu tin được?''

Tôi nghe thấy tiếng cười bí ẩn từ đầu dây bên kia, rồi YAGOO chỉ trả lời một câu, với giọng đầy tự tin của người đã chuẩn bị mọi thứ từ trước: ``Cứ tin tôi đi, đây là một quyết định đúng đắn!''

Tôi lại thở dài trước gợi ý mơ hồ đó, cuối cùng chỉ có thể gật đầu đồng ý, dẫu biết rằng dù có hỏi thêm cũng chỉ nhận được những câu trả lời vòng vo: ``Vâng ạ. Để cháu thử xem sao.''

Tôi cúp máy, ngả lưng xuống đệm một lần nữa, mắt nhìn lên trần nhà với cảm giác bồn chồn khó tả. ``\textit{Mở rộng thêm một tầng à\dots}'' tôi tự nhủ trong lòng, không thể đoán được ông CEO định làm gì với ngôi nhà của mình.

\lang{Etsunan}

Bố tôi quay đầu nhìn tôi, ánh mắt tò mò, rồi đặt remote TV xuống: ``Có chuyện gì vậy, Kuon?''

``Một ý tưởng kỳ quặc từ chú Tanigo. Ông ấy muốn chúng ta mở rộng thêm một tầng cho nhà mình.''

``Mở rộng thêm một tầng?'' bố tôi nhướng mày, sự ngạc nhiên hiện rõ trên khuôn mặt.

``Vâng ạ. Chú ấy nói sẽ có người cần chỗ ở, nhưng không nói là ai.''

Bố tôi gãi cằm, trông đang suy nghĩ sâu xa: ``Nhà mình vốn cũng không rộng rãi gì, mà mở rộng thêm cũng được. Nhưng có một vấn đề\dots''

Ông dừng lại một chút, rồi nhún vai với thái độ thực tế: ``Đầu tiên là tiền đâu ra?''

Tôi chống cằm, lẩm bẩm trong sự băn khoăn: ``Cái đó\dots\ con nghĩ chắc chú Tanigo sẽ hỗ trợ. Nếu không thì chỉ có thể trích một phần lương của con ra thôi.''

Bố tôi bật cười, giọng có chút trêu chọc: ``Bố không nghĩ con có đủ tiền đâu. Mới kiếm được chút ít mà đã định đổ hết vào xây nhà rồi à?''

Tôi thở dài, đưa tay xoa trán: ``Con cũng đâu muốn vậy. Nhưng mà\dots'' tôi ngả đầu ra sau, nhìn lên trần nhà như thể có thể nhìn xuyên qua lớp bê tông để hình dung ra tầng mới: ``Nói gì thì nói, chú Tanigo chưa bao giờ đưa ra yêu cầu vô nghĩa. Có lẽ việc này thực sự quan trọng.''

``Ừ. Bố cũng thấy có lý. Lúc đưa Aloe về nhà mình, chắc chắn cũng phải có lý do gì đó chứ\dots'' bố tôi gật gù, như đang cố ghép các mảnh ghép của câu đố lại với nhau.

``Dù chúng ta còn chẳng biết ông ấy định làm gì,'' tôi thở dài, nở một nụ cười chua chát.

``Rồi còn vô số vấn đề khác sẽ phát sinh\dots\ Xây dựng đâu phải là nghề của bố con mình\dots'' tôi tiếp tục, đặt tay lên trán khi cảm thấy cơn đau đầu âm ỉ bắt đầu xuất hiện.

Khi hai bố con vẫn đang bàn luận đủ thứ vấn đề, cánh cửa trượt phòng đột ngột bật mở. ``Bọn con về rồi đây\~{}!''

Em gái tôi và Aloe bước vào, mỗi người cầm một túi đồ ăn vặt. Nàng Succubus nhanh nhẹn chạy lại gần bàn, rồi nhảy lên đệm với vẻ mặt đầy háo hức.

Có vẻ như Aloe đã nghe lỏm được cuộc trò chuyện của chúng tôi. Cô nàng nở một nụ cười tinh nghịch, đuôi quỷ vẫy vẫy đầy phấn khích: ``Ehehe\~{} Chắc là YAGOO-san muốn chuẩn bị phòng cho em rồi!''

Ngay lập tức, không cần ai bảo, bố tôi ``tặng'' cô nàng một cái cốc nhẹ vào đầu.

*BỐP!*

``Ái da! C-Chú Ryujin! Cháu chỉ đùa thôi mà\dots'' cô nàng phụng phịu, tay xoa chỗ vừa bị gõ, đôi mắt long lanh giả vờ tủi thân.

``Nằm mơ ít thôi, Aloe,'' bố tôi đáp lại, tuy trông có vẻ nghiêm khắc nhưng khóe miệng vẫn nở một nụ cười. ``Chẳng ai rảnh để xây thêm phòng cho cháu đâu.''

Tôi lắc đầu bất lực, góp lời: ``Phòng riêng của em rộng gấp đôi cả phòng của bố và anh cộng lại rồi, còn kêu ca gì nữa.''

``Nhưng mà\dots'' cô nàng cố gắng giữ vẻ buồn rầu, rồi lại bật cười tinh nghịch: ``Ehehe\~{}''

``Thôi đi, Aloe. Đừng có đòi hỏi quá mức,'' tôi thở dài, rồi quay sang bố với vẻ mặt nghiêm túc trở lại: ``Bố, con nghĩ chúng ta nên thử xem sao.''

``Ừ, bố cũng nghĩ vậy.'' bố tôi gật đầu, rồi hỏi em gái tôi: ``Ryuuhou, con thấy sao?''

Em gái tôi nhìn hai người một lúc, rồi ngồi phịch xuống sàn, tò mò hỏi với đôi mắt mở to: ``Khoan đã, hai người đang nói chuyện gì vậy?''

Tôi đập tay lên trán, mệt mỏi giải thích: ``Chú Tanigo muốn nhà mình mở rộng thêm một tầng.''

``Ể? Thêm một tầng nữa á?'' em gái tôi chớp mắt ngạc nhiên, giọng có chút bối rối: ``Nhưng để làm gì vậy?''

Tôi tặc lưỡi, xoa gáy bất lực: ``Không biết. Chú ấy nói sẽ có người đến ở.''

``Ai?''

``Chịu. Đừng hỏi anh. Ông ấy đâu có chịu nói.''

``Tạm bỏ qua chuyện chúng ta xây cho ai đi,'' Aloe lên tiếng, khoanh chân ngồi với vẻ mặt đang suy nghĩ, đuôi quỷ ngừng vẫy như thể đang tập trung: ``nhưng vấn đề là chúng ta sẽ xây thế nào đây?''

``Chắc là\dots\ phải tự xây thôi,'' em gái tôi nói, nhìn bố với vẻ lo lắng.

``Đúng vậy\dots'' bố tôi gật gù, xoa cằm. ``Tiền thì chắc ông sếp của Kuon lo, nhưng làm sao để xây được mới là vấn đề.''

Bố nói chẳng sai. Trong gia đình Hikawa, chỉ có bố và bác Hijaki làm về cơ khí, còn tôi thì làm việc với máy tính.

Nói cách khác, chẳng ai trong gia đình có chút kiến thức nào về kiến trúc hay xây dựng. Ngay cả việc đơn giản là vẽ một bản thiết kế cũng là một thách thức lớn.

Thêm vào đó, chẳng ai trong đại gia đình còn giữ được bản thiết kế gốc của căn nhà này.

``Thế thì\dots\ chúng ta thử vẽ lại từ đầu xem sao?'' em gái tôi đề xuất, giọng nghe hơn như một câu hỏi hơn là lời khẳng định.

``Ừ, thử xem nào,'' tôi gật đầu, lấy giấy bút đặt giữa bàn với chút hy vọng mỏng manh.

Cả bốn người quây quần quanh chiếc bàn gỗ trong phòng khách, bắt đầu buổi họp ``tự phát'' về việc xây tầng mới. ``Được rồi, bắt đầu thảo luận xem tầng mới sẽ được xây ra sao.''

Ngay sau khi tôi vừa dứt lời, Aloe đột ngột đập tay xuống bàn, mắt sáng rực như vừa nghĩ ra ý tưởng tuyệt vời: ``Chúng ta có thể xây một căn phòng có giường treo, cửa sổ toàn kính, trần nhà lấp lánh như bầu trời sao, và cả một bể bơi nhỏ trong phòng nữa!''

Ba người còn lại đồng loạt quay sang nhìn nàng Succubus với vẻ mặt khó hiểu, như thể vừa nghe một câu chuyện cổ tích viễn tưởng.

``\dots Em tưởng nhà mình là khách sạn năm sao à?'' tôi nheo mắt, giọng đầy bất lực. ``Vả lại chúng ta xây cho người khác ở, chứ có phải xây cho em đâu.''

Nghe vậy, Aloe bĩu môi, đuôi quỷ cụp xuống vì thất vọng: ``Đừng dập tắt sáng tạo của em chứ! Phòng hiện đại thì phải hoành tráng chứ! Không thì gọi gì là xây thêm tầng, đúng không?''

``Ít ra cũng phải nghĩ đến cái gì đó thực tế một chút, cháu ạ,'' bố tôi thở dài, với vẻ mặt của người đã trải qua nhiều chuyện.

``Vậy thì xây thành pháo đài đi!'' em gái tôi hào hứng góp ý, mắt sáng rực như đứa trẻ đang tưởng tượng ra cuộc phiêu lưu: ``Làm tường đá, cửa sổ có lỗ châu mai, thêm cả tháp canh nữa!''

``Nhà mình là nhà ở, không phải thành cổ thời Trung Cổ đâu con,'' bố tôi nhăn mặt, lắc đầu bất lực.

``Mà nếu làm vậy thì cả nhà sẽ không chịu được sức nặng đâu,'' tôi bổ sung, cố gắng giữ mọi người bám vào thực tế: ``Tầng dưới là bê tông cốt thép, mà tầng trên làm toàn đá thì có khi cả nhà sập luôn.''

Em gái tôi thở dài, tựa cằm lên bàn với vẻ thất vọng: ``Haizz\dots\ thế thì mất vui quá.''

``Còn nhà gỗ trên cây thì sao?'' Aloe lại hứng chí, mắt lại sáng lên với ý tưởng mới: ``Chúng ta xây phần trên thành căn chòi gỗ, có cầu treo, có bệ nhảy\dots''

``Thế thì khác gì công viên giải trí đâu,'' tôi cắt ngang, xoa trán mệt mỏi. ``Chúng ta xây căn phòng để người ta ở, chứ không phải nơi để chạy nhảy lung tung!''

Nghe tôi mắng, cô nàng có cặp sừng bất đối xứng lại bĩu môi bất mãn, đuôi quỷ vẫy vẫy như đứa trẻ đang phản đối: ``Thế thì hai anh có ý tưởng nào hay hơn không?''

Hai bố con tôi nhìn nhau. Chúng tôi đều biết cần một phương án thực tế, không quá ngớ ngẩn, cũng không khó thực hiện. Tóm lại, không thể để hai cô gái với trí tưởng tượng bay bổng này quyết định mọi thứ.

Bố tôi trầm ngâm một lúc, rồi gõ ngón tay lên mặt bàn với kinh nghiệm của người từng liên quan đến xây dựng: ``Thực ra, nếu nói về kết cấu, tầng nhà bác Hijaki xây là vững chắc nhất.''

``Ý bố là tầng dưới tầng mình đang ở đấy ạ?''

``Chính nó. Hoặc nếu không, lấy tầng nhà Aloe làm chuẩn cũng được.''

Tôi nghe vậy gật gù đồng tình, một tia hy vọng nhen nhóm trong đầu: ``Vâng ạ, con cũng thấy hợp lý. Kết cấu đó vững, chỉ cần cải tiến là có thể thêm một tầng ngay trên mái.''

``Thế thì quyết định vậy đi!'' bố tôi vỗ tay. Cả hai người đàn ông nhà Hikawa đều đồng ý với kế hoạch đơn giản này. Em gái tôi và Aloe nhìn nhau, mặt hơi xịu xuống vì ý tưởng của mình bị bỏ qua, giống như hai đứa trẻ bị lấy mất đồ chơi.

``Nhàm chán quá\dots'' Aloe lẩm bẩm, đuôi quỷ cụp xuống vì thất vọng.

``Nhưng ít ra còn có nhà để ở, còn hơn là có bể bơi rồi cả nhà sập luôn,'' tôi nhún vai, nói theo lẽ thường.

``Và chú nghĩ vị khách kia cũng chẳng cần mấy thứ đó đâu\dots'' bố tôi bổ sung, gật gù với vẻ từng trải.

Cả bốn người cùng gật đầu, thống nhất kế hoạch. Bước tiếp theo là bắt tay vào thực hiện: một công việc sẽ còn vất vả hơn gấp bội. Nhưng ít nhất, chúng ta đã có hướng đi, dù vẫn còn nhiều bỡ ngỡ.

\ct

Vậy là kế hoạch mở rộng căn nhà của gia đình Hikawa chính thức khởi động, với một đội ngũ thi công mà độ ``tin cậy'' thì không thể đỉnh cao hơn được nữa:

Tôi, người được giao trọng trách chỉ huy toàn bộ dự án, nhưng thực chất trong đầu chỉ nghĩ đến cách nào để nhanh chóng xong xuôi mọi việc, sớm được quay về với chiếc điện thoại và trò chơi đang dở dang. 

Bố tôi, dù có kinh nghiệm đọc các bản vẽ, nhưng toàn là bản vẽ cơ khí chứ chẳng hề rành về bản vẽ xây dựng nhà cửa. 

Còn Aloe, với cái danh hiệu tự phong là ``kiến trúc sư succubus'', thì lại sở hữu một trí tưởng tượng bay bổng của một cô gái mười lăm tuổi, chẳng bao giờ chịu bận tâm đến những quy tắc vật lý cơ bản. 

Cuối cùng là em gái tôi, người luôn cặm cụi ghi chép tất cả những sai lầm để rút kinh nghiệm, nhưng đáng tiếc là chẳng có ai chịu nghe theo những bài học đó. Dù vậy, cô bé vẫn kiên nhẫn tiếp tục công việc của mình, như một người thư ký tận tụy, hy vọng những gì ghi chép lại sẽ có lúc phát huy tác dụng.

Bốn người chúng ta cùng nhau leo lên tầng thượng, chính là nơi sẽ được cải tạo thành một căn phòng mới.

Tầng thượng vốn là nơi gia đình Hikawa hay tụ tập mỗi khi có lễ pháo hoa vào dịp năm mới, cũng là chỗ lý tưởng để ngồi hóng mát vào những ngày hè oi bức. Trên đây chẳng có mấy thứ đồ đạc, chỉ có một chiếc bồn bằng inox dùng để chứa nước bơm lên.

``Vậy chúng ta nên bắt đầu từ đâu bây giờ?'' tôi hỏi, mắt nhìn quanh không gian trống trải với một chút băn khoăn chẳng biết bắt tay vào làm sao.

``Trước hết, chúng ta cần kiểm tra lại cấu trúc của tầng này đã,'' bố tôi nêu ý kiến, tay khẽ gõ lên bức tường bê tông. ``Sau đó mới xem xét liệu có thể xây thêm một tầng nữa được hay không.''

``Dạ vâng ạ,'' tôi gật đầu, rồi quay sang hỏi Aloe: ``Em có thể vẽ một bản thiết kế mới không?''

``Dạ có ạ!'' cô nàng Succubus đáp liền, đôi mắt sáng rực như thể vừa được giao một sứ mệnh vô cùng trọng đại. Cô nàng lập tức rút ra một cuốn sổ nhỏ và cây bút, bắt đầu phác thảo bản thiết kế mới, dựa trên những gì hai bố con tôi vừa giải thích.

``Còn bọn chú thì sẽ kiểm tra lại toàn bộ kết cấu,'' bố tôi nói, rồi gọi cô con gái: ``Ryuuhou, con kiểm tra xem các cột trụ và sàn nhà còn vững chắc không nhé.''

``Dạ!'' em gái tôi nhận lệnh, ngay lập tức bắt đầu kiểm tra cấu trúc của tầng thượng với vẻ mặt hết sức nghiêm túc. Cô bé gõ nhẹ vào tường và sàn để lắng nghe âm thanh phản hồi, y hệt như một kỹ sư chuyên nghiệp đang làm việc, dù tôi dám cá em ấy còn chẳng hiểu gì về vật liệu.

Ba người cùng chia nhau công việc: người vẽ, người đo đạc, người tính toán. Mọi việc diễn ra khá suôn sẻ, không gặp trở ngại gì đáng kể, dù cho tất cả chúng ta đều chỉ là những tay mơ với việc xây dựng. Thỉnh thoảng, tôi lại nghe thấy tiếng em gái tôi lẩm bẩm một mình khi con bé nhanh tay ghi chú điều gì đó vào cuốn sổ nhỏ của mình.

``Bây giờ đến lúc tính toán lượng nguyên vật liệu cần dùng rồi,'' bố tôi nói, hai tay khoanh trước ngực khi đứng trước khoảng trống trên tầng thượng.

Tôi cũng gật gù đồng tình, rút điện thoại ra và bắt đầu tập trung: ``Bố nghĩ chúng ta cần bao nhiêu gỗ và xi măng cho công trình này?''

Bố tôi gãi cằm, khuôn mặt trông như đang cố gắng nhớ lại một công thức đã quên mất từ lâu: ``Tầng này phải đủ vững chắc để chịu được toàn bộ lực nặng từ tầng trên đè xuống, nhưng cũng không thể quá nặng kẻo làm yếu đi các tầng bên dưới.''

Tôi ngồi xuống, mắt lia qua không gian của tầng thượng, rồi lại liếc nhìn bảng tính vật liệu từ kho hàng của mẹ. Cứ đang mải mê tính toán, tôi bỗng nhận ra một vấn đề hết sức nghiêm trọng.

``Chúng ta\dots\ không có đủ vật liệu để tự xây đâu.''

Lập tức, cả bốn người đều im lặng. Em gái tôi, lúc này vẫn đang cầm cuốn sổ ghi chép, ngước nhìn tôi với ánh mắt lo lắng. ``Ta có đủ khung thép, cọc và móng, nhưng lại thiếu gạch và xi măng để hoàn thiện công trình,'' tôi tiếp lời, khẽ tặc lưỡi một tiếng vì bất lực.

``Vậy bây giờ phải làm sao đây?'' em gái tôi nhăn mặt, tay đưa lên cằm suy nghĩ.

``Để anh vào kho xem lại có gì tận dụng được không,'' tôi thở dài, rồi dẫn cả nhóm đi xuống tầng trệt.

Chúng ta lần lượt đi qua các kho chứa dọc theo cầu thang, cố gắng lục lọi xem có vật dụng nào có thể dùng cho việc xây dựng không. Em gái tôi, vốn luôn cẩn thận tỉ mỉ, đã cầm theo cây đèn pin và bắt đầu phân loại từng món đồ trong kho một cách có hệ thống, chia thành các nhóm ``có thể sử dụng'', ``có thể tái chế'' và ``vô dụng''.

Sau khi bới tung cả đống đồ đạc trong kho ở các tầng cầu thang, từ những dụng cụ cũ kỹ đến cả mấy thùng sơn đã hết hạn sử dụng, tôi đột nhiên nhớ ra một thứ.

``A! À đúng rồi!'' tôi vỗ tay một cái, làm cả nhóm giật mình. ``Tôi nhớ nhà mình có một chiếc máy xây dựng tự động!``

``Máy xây dựng tự động ạ?'' hai cô gái cùng nhau thốt lên, tỏ ra vô cùng ngạc nhiên vì đây là lần đầu tiên họ nghe đến thứ đó.

Thậm chí bố tôi cũng hơi ngạc nhiên một chút, nhưng rồi ông cũng chợt nhớ ra gia đình mình từng có một cỗ máy như vậy. ``Bố nhớ có vẻ đó là cái máy mà bác Hijaki đã tự chế tạo từ lâu rồi.''

Em gái tôi lại hỏi, giọng vẫn giữ được sự bình tĩnh: ``Nó còn hoạt động tốt không hả bố? Nếu là máy cũ thì chúng ta nên kiểm tra kỹ lưỡng trước khi dùng.''

Nói rồi, bố tôi lôi ra một chiếc hộp kim loại có gắn màn hình nhỏ, cùng với đống bộ phận lỉnh kỉnh của chiếc máy. ``Cái máy bác ấy làm cũng lâu rồi, hình như là hồi bác ấy nâng cấp căn nhà này thì phải, bố nhớ là vậy.''

``Con không thể ngờ nhà mình lại có cả một cỗ máy như thế này\dots'' em gái tôi tròn mắt ngạc nhiên, cô bé cẩn thận kiểm tra từng bộ phận của máy như một chuyên gia đang đánh giá một thiết bị mới lạ.

``Nghe có vẻ hiện đại quá!'' Aloe reo lên, cái đuôi quỷ vẫy vẫy không ngừng vì phấn khích. ``Vậy là chúng ta không cần phải tự tay làm tất cả đúng không anh?''

Tôi ngắm nghía chiếc máy đã phủ đầy bụi một hồi, rồi nhún vai đáp, vừa nói vừa thử bật máy xem có hoạt động được không: ``Ừm\dots\ lý thuyết thì là thế. Nhưng mà vẫn cần phải nhập bản thiết kế vào máy trước đã.''

``Vậy tức là chúng ta phải dùng bản thiết kế mà em vừa vẽ đó, Aloe-chan ạ,'' tôi liếc nhìn cô nàng Succubus với vẻ mặt chẳng mấy tin tưởng.

``Có đây này!'' Aloe reo lên, rồi chậm rãi cúi đầu chỉ vào tờ giấy mà cô nàng vừa đặt trước mặt mọi người.

``\dots Cái gì thế này?'' tôi và bố tôi cùng nhau nhíu mày khi nhìn vào thứ mà Aloe vừa phác thảo.

Trên tờ giấy là một bức tranh đầy đủ màu sắc với vô số chi tiết hết sức kỳ lạ: một tòa tháp hình cây nấm, có những ô cửa sổ hình trái tim, một cái cổng vòm phủ đầy ruy băng lượn sóng, và một bức vẽ nhỏ ở góc trông giống hệt Aloe đang giơ tay cười toe toét.

Em gái tôi nhìn bức tranh, tay đưa lên cằm và đưa ra nhận xét vô cùng khách quan: ``Về mặt thẩm mỹ thì đẹp thật đấy, nhưng về mặt kỹ thuật thì không thể thực hiện được. Sàn nhà không thể nào chịu được trọng lượng của một tòa tháp đá, và cửa sổ hình trái tim còn sẽ gây ra rất nhiều vấn đề về cấu trúc nữa.''

Aloe vẫn hí hửng giải thích, chẳng hề nao núng trước những nhận xét thực tế của em gái tôi: ``Đây là thiết kế của em mà! Em đã dồn hết cả tâm huyết vào đây rồi! Anh chị tưởng tượng xem đi, một căn phòng thật hoành tráng, đẹp lung linh với cửa sổ kính màu lấp lánh, một ban công hình cánh bướm, và đèn chùm lung linh y hệt như trong cung điện hoàng gia!``

Em gái tôi lắc đầu, giọng nói đầy sự thực tế: ``Em thấy có gì đó không ổn lắm\dots\ Em tưởng chị là quỷ mà? Nhưng mà dù là quỷ thì cũng phải tuân theo quy luật vật lý của thế giới này chứ?''

Tôi nhìn bức tranh mà cô nàng vẽ ra, tay tự bóp trán để cố nén tiếng thở dài: ``Aloe-chan\dots\ cái này đâu phải bản thiết kế. Cái này là\dots\ tranh minh họa thôi.''

``Hả?'' cả Aloe và em gái tôi đều ngớ người ra, nhưng em gái tôi nhanh chóng lấy lại được vẻ bình tĩnh.

``Bản thiết kế thì phải có kích thước cụ thể, loại vật liệu rõ ràng, mặt bằng chi tiết\dots\ chứ không phải chỉ là một bức vẽ trông vui mắt là được,'' bố tôi khoanh tay, nói với giọng đầy sự ngán ngẩm. ``Nhìn vào cái này, chú cũng chẳng biết phải xây dựng theo kiểu gì nữa.''

``Ừ thì\dots'' nàng Succubus gãi má cười gượng. ``Ít nhất cháu cũng còn có ý tưởng chứ! Không giống như ai đó suốt ngày chỉ biết ngồi chơi game kia!``

Tôi - kẻ vừa bị nhắc tên - khoanh tay, hỏi với giọng nghiêm nghị như một vị giám khảo khó tính: ``Vậy em có thể nói ra kích thước tổng thể của căn phòng này không? Dài bao nhiêu, rộng bao nhiêu, cao bao nhiêu? Dùng vật liệu gì để xây?''

Trước loạt câu hỏi dồn dập của tôi, Aloe chỉ có thể đứng im lặng, mặt nghệt ra, mắt đảo quanh như đang cố tìm câu trả lời cho một bài toán khó nhằn nhất. ``Cỡ\dots\ ừm\dots\ một căn phòng rộng rãi, đủ để ở là được chứ sao!``'

Bố tôi chỉ có thể cười trừ trước câu trả lời mơ hồ của Aloe, còn tôi thì thở dài ngán ngẩm, vừa vo tròn tờ bản vẽ kia vừa nói: ``Nghiêm túc đi nào. Chúng ta cần một bản thiết kế thực tế hơn thế này nhiều.''

Dù cảm thấy buồn vì ý tưởng của mình không được công nhận, Aloe cũng chỉ có thể ỉu xìu gật đầu: ``Dạ\dots''

Em gái tôi nhẹ nhàng đặt tay lên vai cô nàng, giọng nói đầy sự an ủi: ``Ý tưởng của chị không tệ đâu, nhưng cần phải triển khai nó một cách thực tế hơn thôi. Chúng ta chắc chắn làm được, nhưng phải đi từng bước một chị ạ.''

Nghe vậy, Aloe cũng có vẻ tỉnh táo hơn một chút, nhưng khuôn mặt vẫn còn hơi buồn khi nhận ra bản thiết kế đầy mơ mộng của mình đã bị từ chối một cách nhẹ nhàng nhưng dứt khoát.

\ct

Sau một lúc mày mò, cả bốn người cùng ngồi quây quần, cố gắng phác thảo ra một bản thiết kế sơ đẳng nhất có thể: một phiên bản có đầy đủ thông số, dễ đọc nhưng vẫn bao quát hết các chi tiết cần thiết. Em gái tôi, vốn có tính cẩn thận đến từng li từng tí, đã ghi chép tỉ mỉ từng con số đo đạc lên một tờ giấy riêng, không để lọt bất kỳ sai sót nào.

``Vậy là cuối cùng chúng ta cũng có bản vẽ rồi,'' tôi nói, mắt nhìn vào tờ giấy vừa hoàn thành, lòng dâng lên chút tự hào vì cuối cùng cũng có thứ gì đó trông giống một bản thiết kế đàng hoàng.

``Thế thì chuẩn bị bắt tay vào làm thôi,'' bố tôi lên tiếng, vừa nói vừa bước đến chiếc máy xây dựng tự động mà cả bốn người đã cùng nhau khiêng lên tầng thượng với một quyết tâm mới.

Chúng ta dần dần xếp các vật liệu xây dựng - gỗ, gạch, xi măng - lấy từ kho của mẹ vào máy, từng bước một chuẩn bị cho quá trình thi công. Em gái tôi còn kiểm tra lại từng loại vật liệu một lần nữa, chắc chắn rằng chúng được sắp xếp đúng thứ tự và đủ số lượng như đã tính toán.

``Bắt đầu thôi!''

Tôi nhấn nút khởi động. Chiếc máy lập tức rung nhẹ, phát ra tiếng bíp bíp báo hiệu trước khi chính thức hoạt động. Các cánh tay cơ khí từ từ vươn ra, lần lượt xếp từng tấm gỗ, từng bao xi măng đúng theo vị trí trên bản vẽ. Cả nhóm đều nín thở, mắt dán chặt vào quá trình máy tự động thi công diễn ra trước mắt.

\begin{center}
    \uline{Một tiếng sau.}
\end{center}

Chiếc máy đã hoàn thành việc xây dựng căn phòng mới. Nhưng thứ mà nó tạo ra lại khiến cả bốn người chúng ta trợn tròn mắt, đến mức cùng nhau thốt lên: ``C-Cái quái gì thế này!?''

Căn phòng nghiêng hẳn 45 độ, trông như sắp đổ bất cứ lúc nào. Một phần trần nhà bị lõm xuống, chẳng khác gì vừa bị một thiên thạch nhỏ rơi trúng. Cánh cửa thì cao tận ba mét, trong khi tất cả chúng ta đều chỉ cao hơn một mét bảy một chút. Toàn bộ công trình trông giống như một tác phẩm nghệ thuật trừu tượng được vẽ ra bởi một họa sĩ đang lên cơn sốt.

Em gái tôi xanh mặt, mắt mở toang với vẻ khiếp sợ: ``Em nghĩ\dots\ chúng ta nên dừng hết mọi thứ lại trước khi cả căn nhà của mình biến thành một tác phẩm sắp đặt trừu tượng thật sự.''

Aloe khoanh tay trước ngực, bĩu môi bực bội: ``Vậy bây giờ phải làm sao bây giờ?''

``Cái này\dots\ chú thật sự không bao giờ nghĩ tới được,'' bố tôi toát mồ hôi hột, tay run run chỉ vào căn phòng méo mó trước mặt.

``Nó có vẻ\dots\ không giống với bản thiết kế chúng ta vẽ lắm,'' tôi tặc lưỡi, cố gắng tìm ra một lý do nào đó để bào chữa cho thảm họa vừa ra đời này.

Sau khi nhìn ngắm công trình kỳ quái đó một lúc, cả bốn người đều dần nhận ra một sự thật đáng sợ: ``\dots Nếu cứ để thế này, cả nhà Hikawa có thể sập thật sự.''

Tôi thở dài, vỗ trán suy nghĩ một hồi, rồi với tay lấy điện thoại với vẻ mặt bất lực: ``Có lẽ\dots\ phải nhờ người có kinh nghiệm hơn thôi. Để con gọi bác Hijaki.''

Em gái tôi nghiêng đầu, giọng vẫn giữ được sự bình tĩnh hiếm có: ``Bác Hijaki ạ? Anh hỏi người đã chế tạo ra cái máy này ư? Liệu bác ấy có biết sửa không?''

``Không chắc chắn lắm, nhưng bác ấy cũng là kỹ sư cơ khí giống bố mình. Anh nghĩ bác ấy chắc chắn biết cách sửa chiếc máy xây này.''

Tôi bấm số, chờ đầu dây bên kia nhấc máy. Chỉ vài giây sau, giọng của bác Hijaki vang lên, còn có chút ngái ngủ: ``A-lô? Kuon đấy à? Có chuyện gì vậy?''

``Bác ơi\dots\ bác có rảnh không ạ? Nhà cháu đang có vấn đề to lớn với cái máy xây dựng tự động của bác.''

``Hử? Sao lại thế? Cái máy đó tao cũng đã không dùng từ bao lâu rồi.''

``Dạ, chuyện là thế này\dots\ Căn phòng bọn cháu vừa xây bằng máy của bác, nó bị nghiêng 45 độ, cửa thì cao 3 mét, trần nhà lại bị lõm xuống, và có vẻ như nó sắp sập thật sự.''

Ở đầu dây bên kia, một khoảng im lặng kéo dài. Tôi có thể tưởng tượng ra khuôn mặt của bác Hijaki đang cố gắng xử lý hết những thông tin điên rồ mà tôi vừa đưa ra.

Và rồi\dots

Một tràng cười lớn vang lên. Tiếng cười vang đến mức cả ba người còn lại đứng bên cạnh tôi cũng đều nghe rõ, như thể ông ấy vừa nghe được câu chuyện cười hay nhất trong đời.

Aloe bĩu môi, rõ ràng là không hài lòng, cái đuôi quỷ của cô nàng cũng cụp xuống vì xấu hổ: ``Bác ấy đang cười nhạo chúng ta đúng không?''

``Có vẻ là vậy\dots'' tôi lẩm bẩm, thở dài chán nản và liền giữ điện thoại ra xa tai để tránh bị điếc bởi tiếng cười của ông bác.

Phải mất một lúc bác Hijaki mới bình tĩnh lại được, giọng ông vẫn còn vương chút hài hước: ``Thôi được, tao lên nhà xem sao. Để tao lấy đồ nghề đã.''

``Bọn cháu đang ở trên tầng thượng đấy ạ. Bác lên nhanh nhanh nhé.''

\ct

Chẳng mấy chốc, một tiếng trôi qua, bác Hijaki đã lên trên tầng thượng với một túi dụng cụ đồ sộ đeo sau lưng, dáng vẻ y hệt một kỹ sư chuyên nghiệp chuẩn bị bước vào hiện trường một vụ việc phức tạp. Ngay khi mắt chạm vào căn phòng méo mó đang nằm xiêu vẹo trên sân thượng, ông ấy suýt nữa bật cười ngã ngửa. ``Trời ơi, đây là tác phẩm nghệ thuật cái gì thế hả các cháu!?''

``Thế là mấy đứa định xây cái gì trên sân thượng nhà mình vậy?'' ông ấy vừa cười vừa hỏi, ngón tay chỉ vào công trình kỳ quái như thể đang ngắm nhìn một bức tranh trừu tượng trong bảo tàng.

``Chuyện dài dòng lắm ạ\dots'' tôi thở dài một hơi, rồi lần lượt kể lại mọi sự từ đầu: từ cuộc gọi đột ngột của ông sếp YAGOO, cho đến việc cả nhà phải vội vàng chuẩn bị một tầng mới để đón một vị khách mà đến bây giờ chúng tôi vẫn chưa rõ thân phận. Em gái tôi đứng cạnh bên, lặng lẽ đưa cho bác Hijaki cuốn sổ ghi chép tỉ mỉ của mình để ông có thể nắm được toàn bộ tình hình.

``À, hóa ra là bọn mày định mở rộng nhà thêm một tầng theo yêu cầu của sếp thằng Kuon à?'' bác Hijaki gật gù, dường như đã hiểu rõ cốt truyện.

``Dạ đúng ạ. Nhưng mà đáng tiếc là bọn cháu đã tự ý vẽ bản thiết kế và dùng máy của bác để thử, kết quả thì ra thành một mớ hỗn độn như bác thấy đấy\dots'' tôi đập nhẹ tay lên trán, nhìn cái công trình đáng xấu hổ trước mặt với vẻ ngao ngán tột độ.

``Cái máy này anh cũng đã bỏ quên khá lâu rồi,'' bố tôi tiếp lời, tay vỗ nhẹ vào thân máy cũ kỹ. ``Anh mở ra kiểm tra xem lỗi nằm ở đâu.''

Nghe vậy, bác Hijaki liền mở nắp khoang điều khiển của chiếc máy xây dựng, tỉ mỉ kiểm tra từng bộ phận bên trong, đồng thời thử nghiệm các chi tiết khác với vẻ mặt tập trung của một chuyên gia thực thụ.

Sau khi hí hoáy mày mò một lúc khá lâu, ông ấy rút ra một bảng điều khiển nhỏ, bấm vài nút tùy chỉnh, rồi quyết định cho máy chạy lại một lần nữa để kiểm tra.

Khi chiếc máy bắt đầu vận hành, cố gắng xây dựng lại căn phòng từ đầu, cả năm người chúng tôi cùng nín thở dõi theo từng bước di chuyển của cánh tay cơ khí.

Nhưng tình hình chẳng cải thiện được chút nào. Căn phòng vẫn dựng lên một cách xiêu vẹo, méo mó không khác gì lúc bác Hijaki vừa tới, chưa hề có động tay vào cái máy.

Ông bác lắc đầu ngao ngán, đóng chặt nắp máy lại với vẻ mặt rõ ràng đã tìm ra gốc rễ của vấn đề: ``Không phải lỗi ở máy đâu. Cỗ máy này vẫn hoạt động hoàn toàn ổn. Vấn đề nằm ở dữ liệu mà bọn mày nhập vào đó thôi.''

``Hả!?'' cả bốn người chúng tôi cùng nhau trố mắt ngạc nhiên, không thể tin vào kết luận của bác Hijaki.

Bác Hijaki bật cười phì, xoay màn hình điều khiển của máy về phía chúng tôi, rồi tiếp tục giải thích với giọng đầy thú vị: ``Có người đã nhập sai hoàn toàn các kích thước trong bản thiết kế, nên máy mới xây ra được cái thứ quái dị như thế này.''

``\dots Ai là người đã làm vậy?'' tôi liền liếc mắt sang phía Aloe và em gái tôi. Em gái tôi thì chớp mắt vô tội, giơ hai tay lên trời như muốn nói ``Chắc chắn không phải em!'' với khuôn mặt vô cùng thành thật. Còn Aloe thì lại nhìn lảng sang phía khác, huýt sáo giả vờ như không biết gì, cái đuôi quỷ của cô nàng thì quất qua quất lại một cách rõ ràng là đang giấu giếm điều gì đó.

``\dots Aloe? Liệu có phải em không?''

``Ờm\dots\ có thể là em thật\dots'' nàng Succubus cười gượng, tay đưa lên gãi đầu với vẻ mặt vô cùng xấu hổ. ``Tại em nghĩ cửa nhà cao một chút thì trông sẽ sang trọng, hoành tráng hơn mà\dots''

``Cao tận 3 mét hả?'' bác Hijaki nhíu mày, giọng đầy sự không thể tin nổi. ``Cần gì mà phải làm cái cửa cao đến thế hả?''

Aloe chỉ có thể gãi đầu cười trừ, cái đuôi quỷ cụp xuống vì thất vọng. Bố tôi và tôi thì chỉ có thể lắc đầu ngao ngán trước hành động ngốc nghếch của cô nàng, còn em gái tôi thì chỉ biết cười nhẹ, một nụ cười đầy sự thông cảm với người bạn của mình.

Cuối cùng cũng xác định được rằng vấn đề không nằm ở máy móc, mà chính là ở bản vẽ thiết kế sai lệch. Chúng tôi có thể chỉnh sửa lại các thông số một lần nữa, nhưng\dots

``Chắc chắn rồi, chúng ta phải nhờ sự giúp đỡ từ những người bạn của con ở Khoa Thiết kế Công trình, trường đại học rồi,'' tôi đưa ra quyết định cuối cùng, cảm thấy mình đã đạt đến giới hạn của sự kiên nhẫn sau một loạt sự cố dồn dập.

``Bố nghĩ đó cũng là một ý tưởng rất hay đấy,'' bố tôi đáp lời, gật gù đầy tán thành. ``Bố cũng thấy bản thiết kế tự vẽ của chúng ta vẫn còn nhiều điểm chưa ổn thỏa.''

``Hay là để em tự sửa nó lại bằng sức mạnh đặc trưng của một succubus đi ạ?'' cô nàng sở hữu cặp sừng bất đối xứng lại nhí nhảnh đề xuất, đôi mắt sáng rực lên đầy hy vọng. ``Em đảm bảo là sau khi em sửa, nó sẽ đẹp lung linh lắm!''

``Thôi khỏi cần, cảm ơn ý tưởng của em,'' tôi trả lời thẳng thừng, lập tức dập tắt ngay lập tức niềm hy vọng vừa nhen nhóm trong mắt cô nàng. Aloe chỉ có thể phồng má bực bội nhìn tôi, nhưng cũng chẳng dám cãi lại lời tôi.

\ct

\begin{center}
\uline{Ngày hôm sau\\Phòng 32-H2, Đại học Xây dựng Kawachi.}
\end{center}

Hôm nay tôi có mặt ở trường từ rất sớm, khác hẳn mọi ngày khi tôi chỉ đến để đến lớp. Lần này, mục tiêu của tôi là giải quyết hết mớ rắc rối liên quan đến tầng mới mà chúng ta định xây.

Tôi đi thẳng đến khu vực của Khoa Thiết kế Công trình, nơi tôi đã hẹn trước một vài người bạn học có chuyên môn. Tôi đã nhắn tin sớm để nhờ họ xem xét và đóng góp ý kiến cho bản vẽ nhà của mình - những người này biết rõ điều gì có thể đứng vững và điều gì sẽ sụp đổ, chứ không phải kẻ mộng mơ như tôi và Aloe.

Ngay khi tôi đặt tập tài liệu chứa bản thiết kế lên bàn, cả nhóm sinh viên hào hứng mở ra xem\dots\ và chỉ trong nháy mắt, cả căn phòng im lặng đến khó chịu, như thể họ vừa chứng kiến một vụ va chạm giao thông tàn khốc trước mắt.

``\dots Đây là cái quái gì thế?'' một thành viên trong nhóm thốt lên, mắt mở toang như không thể tin được những gì mình đang nhìn. Lần lượt, những người khác bắt đầu chỉ ra từng lỗi nhỏ to trong bản vẽ, những sai sót thậm chí người mù cũng có thể nhận ra.

``Tôi còn không dám chắc đây là bản vẽ xây nhà, hay chỉ là những nét vẽ nguệch ngoạc của một đứa trẻ lớp một\dots'' một cậu sinh viên khác xoa thái dương, mắt dán chặt vào mớ đường nét méo mó trên giấy, như thể đang cố gắng giải mã một thông điệp gửi từ hành tinh xa xôi.

``Cả đường thẳng đường cong đều lộn xộn hết cả lên\dots'' một người khác chỉ vào những nét vẽ mà Aloe đã vẽ với sự tự tin hệt như một họa sĩ trừu tượng đang triển lãm tác phẩm của mình.

``Khoan đã\dots\ sao ở đây không có một con số nào cả? Chỉ toàn hình vẽ thôi à?'' một người khác hỏi, giọng đầy bối rối không hiểu chuyện gì đang xảy ra.

Tôi chỉ đứng im lặng, để mọi người thoải mái chỉ ra tất cả những sai sót trong bản vẽ. Tôi hiểu quá rõ bản vẽ này mơ hồ đến đâu, và cũng đoán được họ sẽ phẫn nộ như thế nào trước một nửa tác phẩm nửa vời này. Cảm giác lúc đó của tôi thật giống như một bị cáo đang đứng trước bồi thẩm đoàn.

Một cô gái trong nhóm chỉ vào phần vẽ của Aloe, nơi thay vì các tỷ lệ kích thước chính xác, lại xuất hiện một khung cửa hình vòm cực kỳ hoành tráng với hai cánh cửa khổng lồ, đi kèm dòng chữ viết tay ``Lối vào vương giả''. ``Tại sao lại có cái này ở đây? Lấy nhầm bản vẽ rồi phải không?''

``Không đâu, đó đúng là một phần của bản vẽ đấy,'' tôi trả lời, mắt không rời khỏi mớ hỗn độn mà nàng Succubus đã tạo ra, kèm theo một nụ cười chua chát chẳng thể nào vui được.

Sau một hồi nghiên cứu cái bản vẽ lộn xộn này, một cậu sinh viên tiền bối - có lẽ là trưởng nhóm của họ - lắc đầu rồi tuyên bố thẳng thừng: ``Bản vẽ này không thể dùng được.''

``Nghiêm túc đấy. Ai là người đã duyệt cái bản vẽ này vậy?'' một cậu sinh viên đeo kính nhíu mày, giọng điệu hệt như một thẩm phán đang chất vấn kẻ gây rối.

Tôi thở dài, dựa lưng vào ghế với vẻ mặt cam chịu, chuẩn bị đón nhận những ánh mắt phán xét từ mọi người. Tôi khẽ giơ tay, thừa nhận: ``Là bọn tôi tự vẽ đấy.''

Cả nhóm nhìn nhau không thể tin nổi, khi nghe một sinh viên ngành máy tính lại dám tự vẽ bản vẽ xây dựng tồi tệ đến thế. Nó nghe còn phi lý hơn cả câu chuyện cười vô duyên nhất mà họ từng nghe.

Một người khác hỏi, giọng đầy tò mò không dứt: ``Sao lại để cái bản vẽ này đi vào thi công vậy? Ai là người đưa ra ý tưởng thiết kế ban đầu?''

Tôi nhún vai, thú nhận thật thà với giọng đầy bất lực: ``Gia đình tôi muốn mở rộng nhà thêm một tầng, nhưng chẳng ai còn giữ bản thiết kế gốc của căn nhà từ lâu rồi.''

Lần nữa, cả nhóm sinh viên im lặng nhìn nhau, như thể họ vừa thấy một kẻ không có chút kinh nghiệm nào lại đòi tự sửa chiếc xe hơi của mình, y hệt một đứa trẻ muốn đóng vai thợ máy.

Một thành viên trong nhóm thở dài, lắc đầu với vẻ mặt của một chuyên gia đang tìm cách cứu vãn một tình huống đã gần như hỏng hết: ``Được rồi\dots\ Với tình trạng này thì không thể sửa qua loa được đâu. Nếu muốn mở rộng mà vẫn đảm bảo an toàn kết cấu, tốt nhất là bọn tôi phải đo đạc lại toàn bộ căn nhà từ đầu.''

Một cô gái khác gật đầu lia lịa, giọng điệu chắc nịch của một kỹ sư giàu kinh nghiệm: ``Đúng vậy. Không có bản vẽ gốc, chúng tôi phải kiểm tra lại tất cả: kích thước từng phòng, vị trí các bức tường chịu lực, độ sâu và độ bền của móng nhà, rồi mới có thể vẽ lại bản thiết kế đúng chuẩn được.''

Tôi lại một lần nữa thở dài, tựa đầu vào lòng bàn tay, khẽ chẹp miệng như thể đã lường trước được kết quả này từ lâu: ``Đành vậy thôi, còn cách nào khác đâu.''

``Nhắc trước là công việc này khá vất vả đấy nhé,'' một người trong nhóm cười khẩy, như thể đang cảnh báo tôi về một cuộc hành trình dài và đầy mệt nhọc phía trước.

``Rồi, tôi biết rồi\dots''

Cuối cùng, sau một hồi bàn bạc đi đi lại lại, nhóm sinh viên thống nhất sẽ đến nhà Hikawa để thực hiện toàn bộ công việc đo đạc trong vòng một tuần tới. Tôi chỉ có thể gật đầu đồng ý, chấp nhận rằng kế hoạch mở rộng nhà này đang phức tạp và tốn nhiều công sức hơn tôi từng nghĩ.

\ct

Chiều cùng ngày, cả nhóm sinh viên từ Khoa Thiết kế Công trình có mặt tại nhà Hikawa, mang theo đủ thứ thiết bị đo đạc chuyên dụng trông như những cỗ máy khoa học viễn tưởng. Bố tôi và tôi đứng đón họ ở cửa, vừa mừng rỡ vì sắp có chuyên gia giúp đỡ, lại vừa hơi ngượng nghịu, vì sau cùng, chúng tôi đang nhờ vả họ làm việc miễn phí mà.

``Bắt đầu thôi. Đầu tiên chúng ta sẽ đo kích thước từng tầng, sau đó mới lên kế hoạch chi tiết cho tầng mới,'' một người trong nhóm vỗ tay ra lệnh, hệt như một vị chỉ huy đang khởi động chiến dịch quan trọng.

Một nhóm cầm thước laser và các máy đo chuyên dụng, bắt đầu kiểm tra từng chi tiết: chiều cao trần nhà, vị trí chính xác của các bức tường chịu lực, cả cách bố trí cầu thang giữa các tầng; họ làm việc tỉ mỉ không kém những người thợ lành nghề có nhiều năm kinh nghiệm. Trong khi đó, nhóm còn lại ngồi ở phòng khách, bắt tay vào vẽ lại bản thiết kế chuẩn mực, chứ không phải một bức tranh đầy tưởng tượng nguệch ngoạc như tác phẩm của Aloe.

Với bố tôi và tôi, chúng tôi vẫn giữ nguyên ý tưởng ban đầu: lấy căn phòng của bác Hijaki ở tầng bốn làm mốc. Đó là phòng được thiết kế bài bản nhất, và cũng dễ dàng tái sử dụng sau này nếu có nhu cầu. Các sinh viên cũng đồng tình rằng đây là phương án an toàn nhất, khiến cả hai chúng tôi đều thở phào nhẹ nhõm.

Một tuần trôi qua yên bình đến lạ, một điều cực kỳ hiếm gặp trong căn nhà Hikawa luôn ồn ào này.

Tại sao lại thế?

Đơn giản là vì em gái tôi đi học suốt ngày, còn Aloe thì vừa phải đến lớp, vừa phải đi làm thêm ở quán cà phê gần nhà.

Điều đó có nghĩa là chẳng còn ai có thể gây rối hay phá đám công việc, để tôi và bố tôi có thể đẩy nhanh tiến độ mà không bị gián đoạn. Mọi thứ diễn ra hết sức chuyên nghiệp: bản thiết kế dần hoàn thiện, tất cả các thông số đều chính xác, và kết cấu của tầng mới cũng đã được tính toán kỹ lưỡng đến từng chi tiết nhỏ.

Cuối cùng, vào ngày cuối cùng của tuần, Aloe trở về nhà với vẻ mặt đầy tự tin vô đối. Cô nàng bước vào phòng khách với dáng điệu của một tiểu quý tộc mới trở về dinh thự, y hệt như hồi mới chuyển vào nhà Hikawa, nhìn xuống nhóm người đang hoàn thiện những nét cuối cùng của bản thiết kế.

``Ta đã trở về rồi đây, lũ hạ dân!'' cô nàng vênh mặt, cái đuôi quỷ vẫy vẫy đầy kiêu hãnh. ``Chắc trong suốt một tuần qua, các ngươi nhớ đến ta phát điên lên rồi đúng không?''

Cả căn phòng chìm vào im lặng trong vài giây. Không ai thèm để ý đến sự xuất hiện của cô nàng Succubus, như thể cô chỉ là một bóng ma vô tình đi ngang qua.

Em gái tôi vừa bước ra từ bếp, cầm trên tay ly nước, chỉ nhún vai với vẻ mặt cười trộm: ``Em nghĩ chẳng ai để ý đâu, Aloe-oneechan. Đừng làm phiền họ đang làm việc nữa.''

Tôi thậm chí còn không thèm ngẩng đầu khỏi đống tài liệu, mắt vẫn dán chặt vào các con số trên bản vẽ: ``Ai rảnh mà nhớ chứ? Bọn anh còn đang bận chết đi được đây.''

Aloe thoáng khựng lại, cái đuôi quỷ cụp xuống một chút vì thất vọng. Nhưng cô nàng nhanh chóng lấy lại phong thái kiêu ngạo: ``Hừm! Thật đáng thương cho các ngươi, phải vất vả thế này mà không có ta ở đây chỉ dẫn. Nhưng thôi, ta rộng lượng tha thứ, để ta xem các ngươi đã làm được gì rồi.''

Cô nàng chắp tay sau lưng, đi vòng quanh bàn như một bà chủ đang kiểm tra công việc của đám đầy tớ, rồi dừng lại đúng trước bản thiết kế cuối cùng.

Đó là một bản vẽ chi chít đủ thứ thông số: diện tích, kích thước, loại vật liệu, tất cả mọi chi tiết cần thiết đều có đủ. Một bản thiết kế thực sự có thể đưa vào sử dụng, một bản vẽ của người chuyên nghiệp.

``Ồ, cuối cùng cũng ra hồn được một cái bản vẽ. Nhưng ta cá là nếu có ta ở đây từ đầu, các ngươi đã xong sớm từ ba ngày trước rồi.''

Lập tức, một thành viên trong nhóm giơ tay cốc thẳng vào đầu cô nàng một cái, động tác nhanh như chớp. ``Đau!''

``Mơ mộng ít thôi, nhóc con,'' cậu bạn ấy nói. Bố tôi và tôi chỉ liếc nhìn cô gái đang xoa đầu xù xì kia mà ngán ngẩm, còn em gái tôi thì cười nhạt: ``Đã bảo rồi mà có nghe đâu.''

Cô nàng Succubus trừng mắt đầy phẫn nộ, cái đuôi quỷ dựng đứng lên hệt như một con mèo bị chọc giận, lớn tiếng chất vấn: ``Ngươi dám đánh ta ư!? Ta sẽ nhớ mặt ngươi, đừng có mà hối hận!''

Cậu sinh viên vừa gõ đầu cô nàng chỉ cười khẩy một tiếng, chẳng thèm đáp lại, như thể đã quá quen với những phản ứng thái quá của Aloe.

Tôi khoanh tay, nhìn chằm chằm vào cô nàng với ánh mắt của một người đã quá mệt mỏi với đủ thứ trò đùa ngốc nghếch: ``Nói lại cho rõ nhé, phòng riêng của em rộng gấp đôi cả phòng tôi và phòng bố tôi cộng lại, còn đòi hỏi gì nữa nào?''

Aloe cứng họng ngay lập tức. Cô nàng há miệng định cãi lại, nhưng khi nhớ lại căn phòng của mình ở tầng hai, thực sự nó rộng hơn cả hai phòng của tôi và bố tôi cộng lại, cô nàng chỉ biết cúi đầu nhìn xuống sàn với vẻ mặt bối rối.

Cô nàng tặc lưỡi, bực bội đáp lại mà giọng đã yếu đi nhiều: ``\dots .Thật là tức quá! Em không muốn nói chuyện với các anh nữa!'' Aloe hất tóc, cố gắng giữ vẻ ngoài kiêu ngạo nhưng thực chất đã hết lý lẽ để cãi. Cô nàng chỉ có thể ngồi xuống chiếc đệm gần đó, khoanh tay nhìn mọi người làm việc với khuôn mặt vẫn còn cau có.

Em gái tôi tủm tỉm cười, giọng đầy vẻ hài hước: ``Biết điều như vậy từ đầu có phải ai cũng nhẹ nhõm không.''

Dù sao đi nữa, kế hoạch mở rộng căn nhà của gia đình Hikawa cuối cùng cũng đã hoàn thành giai đoạn chuẩn bị. Bản thiết kế đã sẵn sàng, nguyên vật liệu cũng được tính toán đủ đầy, và đội ngũ thi công - dù chủ yếu là những người không chuyên - cũng đã có một lộ trình rõ ràng để đi tới.

\ct

Sau khi chốt xong bản thiết kế cuối cùng, cả nhóm lại cùng nhau trèo lên tầng thượng một lần nữa. Lần này thì chẳng còn những sai sót ngốc nghếch nào xảy ra nữa: bản vẽ đã chính xác từng milimét, còn chiếc máy xây dựng tự động thì đã được Ryujin bảo trì tận tình, và Kuon cũng đã nâng cấp thêm vài tính năng hiện đại để nó hoạt động mượt mà hơn.

``Bắt đầu thôi nào, mọi người sẵn sàng chưa!''

Chiếc máy được khởi động, lập tức phát ra một chuỗi tiếng động cơ khí trầm ấm quen thuộc. Màn hình điều khiển quét nhanh qua các thông số trên bản thiết kế, xác nhận tất cả dữ liệu chính xác trước khi tiến hành thi công chính thức. Cả đám đứng đó nín thở, mắt không rời khỏi từng chuyển động của máy.

``Mọi thứ đã ổn hết rồi,'' tôi thở sâu một hơi, rồi tay ấn mạnh vào nút khởi động.

Lần này, máy vận hành trơn tru đến khó tin. Căn phòng mới được xây dựng nhanh chóng và chính xác tuyệt đối, chẳng còn lỗi vặt như lần trước nữa. Những cánh tay cơ khí di chuyển uyển chuyển, đặt từng viên gạch, từng tấm vật liệu đúng vào vị trí đã định, độ chính xác cao đến từng milimét.

Không chỉ dựng xong căn phòng đúng như bản vẽ, máy còn tự động xử lý luôn một khoảng trống bên cạnh để làm lối đi cho cầu thang dẫn lên tầng trên. Các tia laser định vị vạch ra đường cắt hoàn hảo, rồi cánh tay máy nhẹ nhàng gỡ bỏ một phần trần nhà cũ, tạo ra một khoảng trống vừa vặn để lắp cầu thang sau này.

Tiếp theo, các vật liệu mới lần lượt được lắp ghép vào đúng chỗ: những thanh dầm thép được hàn cố định để gia cố kết cấu, lớp sàn cũ của tầng thượng dần được thay thế bằng nền móng vững chắc hơn cho tầng mới. Tường và trần nhà được dựng lên từng lớp, từ khung sườn chính cho đến lớp cách nhiệt, cuối cùng là lớp sơn hoàn thiện. Cánh cửa lần này có kích thước chuẩn, chẳng còn là cái cửa khổng lồ cao ba mét như lần trước nữa.

Những viên gạch lát sàn, tấm ván gỗ đều được lắp đặt tỉ mỉ, từng chi tiết nhỏ đều được kiểm tra kỹ lưỡng trước khi hoàn thiện. Ánh sáng lấp lánh từ máy hàn kết hợp với tín hiệu của các cảm biến hồng ngoại tạo nên một khung cảnh công nghệ đầy ấn tượng.

Chỉ sau khoảng một tiếng rưỡi thi công tự động, toàn bộ quá trình xây dựng đã hoàn hảo kết thúc. Không còn những góc cạnh méo mó hay những lỗi kỹ thuật kỳ quái nào nữa.

``Và đây là thành quả của chúng ta!'' mọi người cùng vỗ tay reo hò, nhìn căn phòng mới tinh trước mặt với vẻ phấn khích.

Một căn phòng mới tinh tươm, vững chãi, hoàn toàn đúng theo thiết kế đã vẽ. Những bức tường thẳng tắp, cửa sổ đặt đúng vị trí, sàn nhà phẳng phiu và chắc chắn. Một không gian sống thực thụ, hoàn toàn có thể ở được ngay.

Cả nhóm cùng nhau thở phào nhẹ nhõm, như thể vừa trút được một gánh nặng lớn khỏi vai sau bao ngày vất vả.

Bố tôi khoanh tay đứng nhìn, gật gù hài lòng: ``Đúng là có kỹ thuật chuyên nghiệp thì khác hẳn.''

Bác Hijaki bật cười, giọng đầy vẻ hóm hỉnh: ``Thế mới biết, chịu khó nhờ chuyên gia ngay từ đầu có phải đỡ khổ không. Lần sau đừng có tự ý làm nữa đấy.''

Tôi vươn vai thư giãn, cảm nhận sự nhẹ nhõm lan tỏa sau bao ngày lo lắng: ``Cũng may là cuối cùng mọi thứ cũng xong xuôi đâu vào đấy.''

Em gái tôi cười tươi, đôi mắt cô bé sáng rực lên niềm tự hào: ``Vậy là cuối cùng chúng ta cũng có thêm một tầng mới rồi nhỉ.''

Riêng Aloe là người duy nhất trông có vẻ không mấy vui vẻ. Cô nàng khoanh tay, nhăn mặt tỏ ra bất mãn: ``Hừ! Nếu mà làm theo ý tưởng của em thì căn phòng này còn hoành tráng hơn nhiều, chứ không phải nhàm chán thế này!''

Tôi từ từ quay sang nhìn cô nàng, giọng đầy vẻ thách thức: ``\dots Ồ thế sao? Nếu thế thì em định làm kiểu gì nào?''

Nàng Succubus lập tức giơ tay lên hào hứng như một nhà thiết kế thời thượng, bắt đầu mô tả với đôi mắt sáng rực: ``Em đã hình dung ra một căn phòng phong cách lâu đài ma thuật! Trần nhà cao gấp đôi bình thường, có đèn chùm pha lê lung linh, cửa sổ kính màu cầu vồng, rèm nhung đỏ rực rỡ, thậm chí còn có cả một lối đi bí mật dẫn đến thư viện riêng cơ!''

Bác Hijaki nghe vậy suýt nữa sặc nước, bố tôi thì chỉ biết xoa đầu bất lực. Tôi cũng chỉ biết lắc đầu ngán ngẩm.

``Được rồi, mơ mộng ít thôi em ơi,'' tôi dập tắt ngay tắp lự tưởng tượng của cô nàng, giọng thực tế hết mức. ``Chúng ta chỉ cần một căn phòng đủ để ở thôi, chứ không phải xây lâu đài cho quỷ vương đâu.''

Aloe hậm hực, bĩu môi trông như một đứa trẻ bị tước mất cây kẹo yêu thích: ``Mấy người thật không có chút óc thẩm mỹ nào cả! Chỉ biết ở nhàm chán!''

Em gái tôi nhún vai tủm tỉm cười: ``Ít nhất thì bây giờ cả căn nhà vẫn đứng vững là được rồi. Aloe-san sắp có bạn mới chuyển vào ở chung rồi đấy, vui lên đi.''

``\dots Cứ nói mãi chuyện đó đi,'' Aloe gằn giọng giả vờ bực mình, nhưng chẳng thể giấu được vẻ hồ hởi hiện lên trên mặt, cái đuôi quỷ của cô cũng bắt đầu vẫy vẫy không ngừng vì tò mò.

Cả nhóm cùng quay lại nhìn căn phòng mới vừa hoàn thành. Không gian bên trong vẫn còn trống trơn, nhưng nó đã sẵn sàng để đón một người mới chuyển đến. Một chương mới, một tương lai khác đang chờ đợi phía trước.

Một cơn gió nhẹ lùa qua cửa sổ, làm lay động tấm rèm vừa mới treo lên. Ánh nắng chiều vàng ươm hắt vào căn phòng trống, như một lời dự báo tốt lành cho những sự thay đổi sắp tới.

Dù không ai nói ra miệng, nhưng tất cả mọi người đều thầm hiểu: căn phòng này sẽ không trống mãi đâu.

Một vị khách họ Rồng nào đó sẽ sớm trở thành chủ nhân của căn phòng mới toanh này.

Vấn đề chỉ là, ngày nào cô ấy sẽ tới thôi.

\end{document}